\Needspace{30\baselineskip}
\section{Gradients against positive currents}
\label{sec:bt-current-gradients}
The following local chain rule also applies to the cutoff arguments on compact manifolds. Here $\mathrm d^{\mathrm c}=(\partial-\bar\partial)/(4\pi\mathrm i)$, so $\mathrm d\mathrm d^{\mathrm c}=\ddc$.

\Needspace{23\baselineskip}
\begin{lemma}\label{lem:bounded-psh-current-chain-rule}
Let $\Theta$ be a closed positive current of bidimension $(1,1)$ on a coordinate domain, and let $v$ be locally bounded psh. The current $\mathrm d(v\Theta)$ has measure coefficients, denoted $\mathrm dv\wedge\Theta$. There is a positive gradient measure
\begin{equation}\label{eq:bounded-current-gradient-measure}
 E_\Theta(v)=\mathrm dv\wedge\mathrm d^{\mathrm c}v\wedge\Theta
       =\tfrac12\ddc(v^2\Theta)-v\,\ddc v\wedge\Theta.
\end{equation}
Smooth decreasing approximations have gradients converging strongly in the coefficient-matrix energy seminorm on compact subsets. For every smooth real $f$, the bounded products satisfy
\begin{equation}\label{eq:bounded-current-chain-rule}
 \mathrm d(f(v)\Theta)=f'(v)\,\mathrm dv\wedge\Theta.
\end{equation}
For a compactly supported smooth real one-form $\eta$, a cutoff $\tau\geq0$ equal to one near its support, and a bounded Borel function $a$, there is a constant $C_\eta$ such that
\begin{equation}\label{eq:current-gradient-cauchy-schwarz}
 \left|\int a\eta\wedge\mathrm dv\wedge\Theta\right|^2
 \leq C_\eta\left(\int\tau E_\Theta(v)\right)
                 \left(\int_{\operatorname{supp}\eta}|a|^2\,
                                         \Theta\wedge\beta\right).
\end{equation}
\end{lemma}
\begin{proof}
All assertions are local near a compact test support. Add a constant to make $v$ and its smooth radial approximations $v_l\downarrow v$ nonnegative and uniformly bounded there. For a smooth real function $w$, the differential product rule gives
\[
 \mathrm dw\wedge\mathrm d^{\mathrm c}w\wedge\Theta
       =\tfrac12\ddc(w^2\Theta)-w\ddc w\wedge\Theta.
\]
The measure on the left is positive. Applied to $w=v_l-v_m$ and tested with $\tau$, this identity gives energy tending to zero as $l,m\to\infty$. Indeed \cref{cor:bt-polynomial-weights} makes the squared-weight term tend to zero after integration by parts against $\ddc\tau$, and each of the four expanded Hessian terms tends to the same integral with weight $v$ and Hessian $\ddc v$. Their signs cancel. For two independent indices the convergence is uniform on tails: otherwise select violating pairs with both indices strictly increasing and apply that same decreasing-continuity theorem.

Here is the meaning of the energy seminorm, including when $\Theta$ is singular. Its positive coefficient matrix has bounded Radon--Nikodym entries relative to $\mu=\Theta\wedge\beta$, by \cref{thm:positive-current-order-zero}. Write that measurable Hermitian matrix as $H$. For this bidimension, its quadratic form is the density of $\mathrm i\xi\wedge\bar\xi\wedge\Theta$ for a constant $(1,0)$-form $\xi$. Positivity for a countable dense family of $\xi$ makes $H$ positive semidefinite almost everywhere. In coordinates the gradient energy is the squared $L^2(\mu)$ norm of $H^{1/2}\partial w$, with the fixed normalization absorbed into $H$. The preceding estimate makes these vectors for $v_l$ Cauchy on compact subsets. Their strong limit defines the gradient modulo the zero seminorm. Multiplication by a bounded measurable scalar respects that quotient. Pointwise matrix Cauchy--Schwarz bounds the derivative paired with $a\eta$ by the two norms in \eqref{eq:current-gradient-cauchy-schwarz}. It also bounds every derivative coefficient in variation. Distributional convergence of $v_l\Theta$ to $v\Theta$, by dominated convergence, identifies that measure derivative with $\mathrm d(v\Theta)$.

The squared norms converge in $L^1(\mu)$, so their limit is a positive measure. Polynomial weighted continuity identifies its distribution with the right side of \eqref{eq:bounded-current-gradient-measure}. This also proves independence of the smooth approximations. Finally $f(v_l)\Theta\to f(v)\Theta$ by dominated convergence. The products $f'(v_l)$ times the gradients converge strongly: subtract the limiting gradient first, then use dominated convergence on its squared energy with $f'(v_l)\to f'(v)$. Passing through the smooth chain rule proves \eqref{eq:bounded-current-chain-rule}. A constant shift only translates the argument of $f$ and does not change the gradients, completing the proof.
\end{proof}

\index{gradient measure!positive currents}\index{chain rule!bounded psh potentials}
